\section*{序論}
11月1日朝までの生成AI関連Xでは、直近に公開された大規模モデルの評価と、AI学習データの利用制御が並行して話題になった。Mistral Large 4のオープンウェイト公開とGemini 4 Carbonの限定提供について、実際のコーディング能力、サイバー防御への適用、利用条件が比較されている。いずれも主要な発表自体は今回の観測窓より前であり、窓内では利用報告や反応が中心となった。

同時に、OpenAI、Google、Microsoft、Anthropicが共通のrobots.txt信号を尊重するという発表が共有され、出版社やサイト運営者が学習利用を制御する仕組みに関心が集まった。Anthropicの自動レッドチーミング研究やGoogle DeepMindの英国自治体向けPlanAssistは、モデル能力を安全性評価や行政業務へ接続する動きとして注目された。

本号はGrokの07:00収集結果に含まれる8トピックを掲載順に再構成した。未確認の性能比較、実装予定、コミュニティによるモデル一覧は投稿上の主張として扱い、一次発表と観測窓内の反応を区別する。独立したWeb検証は行っていない。
